% !TeX program = xelatex
\documentclass[11pt,a4paper]{article}
\usepackage[margin=23mm]{geometry}
\usepackage{fontspec}
\setmainfont{Latin Modern Roman}
\usepackage{amsmath,amssymb,mathtools}
\usepackage{xcolor,enumitem,booktabs,longtable,array,xurl,hyperref}
\definecolor{muted}{HTML}{53636C}
\hypersetup{colorlinks=true,urlcolor=blue,linkcolor=blue}
\setlength{\parindent}{0pt}
\setlength{\parskip}{5pt}
\newcommand{\R}{\mathbb R}
\newcommand{\E}{\mathbb E}
\newcommand{\N}{\mathbb N}
\newcommand{\ind}{\mathbf 1}
\newcommand{\Var}{\operatorname{Var}}
\newcommand{\Cov}{\operatorname{Cov}}
\newcommand{\supp}{\operatorname{supp}}
\DeclareMathOperator*{\argmin}{arg\,min}
\DeclareMathOperator*{\argmax}{arg\,max}
\newcommand{\entrymeta}[3]{{\sffamily\small\color{muted}\textbf{\texttt{#1}}\quad|\quad #2\par #3\par}}
\newcommand{\Problemref}[1]{\texttt{#1}}
\newcommand{\JELref}[2]{\texttt{#2} (JEL \texttt{#1})}
\begin{document}
\section*{Minimax policy regret under partial causal identification}
\textbf{Problem ID:} \texttt{C44-87}\quad \textbf{JEL:} \texttt{C44}\par
% BEGIN ECONOMICS STATEMENT
\label{problem:OP-0511}
\label{beyond:CI-32}
\entrymeta{CI-32}{Econometric theory and causal inference}{Research agenda formalization}
\subsection*{Definitions and assumptions}
Observe $n$ independent copies of $(X,W,Y)$ with common law $P$, $X\in[0,1]^d$, $W\in\{0,1\}$, $Y\in[-1,1]$, and observed propensity $e_P(X)\in[\varepsilon,1-\varepsilon]$, $P_X$-almost surely, where $e_P(x)=P(W=1\mid X=x)$ and $0<\varepsilon<1/2$ is fixed. Do not assume exchangeability given $X$. For a fixed finite $\Gamma\geq1$, let $\mathcal Q_\Gamma(P)$ comprise full causal laws $Q$ inducing the observed law $P$, satisfying consistency $Y=Y(W)$, no interference, and $(Y(0),Y(1))\perp W\mid(X,U)$ for some latent $U$. Define $e_Q(x,u)=Q(W=1\mid X=x,U=u)$ and require, $Q_{X,U}$-almost surely,
\[
 \Gamma^{-1}\leq
 \frac{e_Q(X,U)/(1-e_Q(X,U))}{e_P(X)/(1-e_P(X))}
 \leq\Gamma.
\]
Assume this set is nonempty. Specify a nonempty, pointwise separable, finitely encoded policy class $\Pi$ of measurable maps $[0,1]^d\to\{0,1\}$, of Vapnik--Chervonenkis dimension $v$, with a stated evaluation and optimization oracle. Pointwise separability means that a countable subclass has, for each $\pi\in\Pi$, a sequence converging to $\pi$ at every covariate point. Let $Z=(Y+1)/2$. Let $\mathcal P$ comprise the observed laws satisfying the stated support, overlap and causal-compatibility requirements, with Lebesgue density $f_X\in\mathcal H_d^s(L)$ and, for each $w\in\{0,1\}$, a conditional Lebesgue density of $Z$ given $(X,W=w)$ admitting a jointly Hölder version $f_{Z\mid X,W=w}\in\mathcal H_{d+1}^s(L)$ on $[0,1]^{d+1}$, for $d\in\N$, $L>0$ and $s>0$. The conditional-density requirement is imposed almost everywhere under the law of $X$ given $W=w$; the Hölder version on the entire cube is an extension requirement outside that support. Restrict attention to parameter choices for which $\mathcal P$ is nonempty. No positive lower density bound or propensity smoothness is imposed. Write $V_Q(\pi)=\E_Q[Y(\pi(X))]$ for the full-law policy value.
\subsection*{Formal question}
Define the robust regret objective
\[
 R_\Gamma(P,\pi)=\sup_{Q\in\mathcal Q_\Gamma(P)}
 \left\{\sup_{\pi'\in\Pi}V_Q(\pi')-V_Q(\pi)\right\}.
\]
Determine the minimax excess robust regret
\[
 \inf_{\widehat\pi:\,\widehat\pi\in\Pi}\sup_{P\in\mathcal P}
 \E_{P^{\otimes n}}\left[R_\Gamma(P,\widehat\pi)-
       \inf_{\pi\in\Pi}R_\Gamma(P,\pi)\right],
\]
and construct a statistically and computationally feasible learner attaining the bound. Characterize dependence on sensitivity strength, overlap, smoothness, and policy complexity. For every rate assertion, state which of $d,L,s,v,\Gamma,\varepsilon$ and the policy-class description are fixed as $n$ grows, and give the parameter ranges and constant dependence for any asserted uniform bound. To formulate dependence of adverse laws on policies without assuming attainment, put $D_Q(\pi,\pi')=V_Q(\pi')-V_Q(\pi)$ and, for $\delta\geq0$, define
\[
 A_\delta(P;\pi,\pi')=
 \left\{Q\in\mathcal Q_\Gamma(P):
 D_Q(\pi,\pi')\geq
 \sup_{\widetilde Q\in\mathcal Q_\Gamma(P)}D_{\widetilde Q}(\pi,\pi')-\delta\right\}.
\]
Characterize when a collection of competing policy pairs has a common exact maximizing law, meaning a nonempty intersection of its $A_0$ sets, and when it has a common $\delta$-maximizing law for each $\delta>0$. Distinguish these assertions from the mere selection of different maximizers from nonunique sets. In every pairwise comparison the same $Q$ governs both policy values; the displayed robust-regret objective and its supremum order remain the target.
\subsection*{Scope and status}
The sensitivity class is an editorial choice, not an assumption supplied by the agenda. Partial identification does not itself prevent learning a robust decision, but its benchmark differs from the optimum under an unknown single causal law. Suprema and infima do not presume existence of optimizers. Established special cases of robust policy learning are not certified as open by this entry.
\subsection*{References for the source of the problem}
Carlos Cinelli, Avi Feller, Guido Imbens, Edward Kennedy, Sara Magliacane, and Jose Zubizarreta, \emph{Challenges in Statistics: A Dozen Challenges in Causality and Causal Inference}, arXiv:2508.17099v1 [stat.ME], 23 August 2025, \S3.3.3, ``New optimization goals,'' including minimal worst-case regret under weakened identification. \url{https://arxiv.org/html/2508.17099v1#S3.SS3.SSS3}.

\subsection*{Common conventions: Source mathematical conventions}

\textbf{Well-defined objects.}
All probability spaces and measurable variables are specified or inherited
from a local statistical model. Expectations used as finite targets require
the stated integrability. A decision rule or estimator is measurable with
respect to its available information; randomized rules may use an auxiliary
independent random variable. Norms without a subscript are Euclidean in
finite-dimensional spaces. An optimizer is requested only when attainment
is assured; otherwise the question concerns the optimal value and
$\varepsilon$-optimal admissible rules for every $\varepsilon>0$.

\textbf{Identification.}
A structural model $m\in\mathcal M$ implies an observable law
$P_m^{\rm obs}$ and target $\theta(m)$. The target is point identified on
$\mathcal M$ if
\[
 P_{m_1}^{\rm obs}=P_{m_2}^{\rm obs}
 \quad\Longrightarrow\quad\theta(m_1)=\theta(m_2)
 \qquad(m_1,m_2\in\mathcal M).
\]
When this fails, the sharp identified set is
\[
 \Theta(P;\mathcal M)=\{\theta(m):m\in\mathcal M,
                                  P_m^{\rm obs}=P\}.
\]
Specifying an intervention or estimand does not establish identification.
Positivity, exclusion, causal sufficiency, measurement restrictions, and
transport assumptions are substantive local inputs, never automatic
consequences of notation.

\textbf{Minimax risk and nontrivial inference.}
For a statistical experiment with data law $P^{(n)}$, target $\theta(P)$,
and nonnegative loss $\ell$, define
\[
 \mathcal R_n(\mathcal P,\ell)
   =\inf_{\widehat\theta_n}\sup_{P\in\mathcal P}
      \E_{P^{(n)}}\ell(\widehat\theta_n,\theta(P)).
\]
The infimum is over measurable estimators. A sharp rate is a sequence
$r_n$ for which $0<c\leq C<\infty$, independent of $n$, satisfies
$c r_n\leq\mathcal R_n\leq C r_n$ for all sufficiently large $n$.
Squared-error parametric risk is of order $n^{-1}$; the corresponding
root-mean-squared error is of order $n^{-1/2}$. These different conventions
are distinguished locally. A uniform asymptotic confidence region $C_n$
must satisfy
\[
 \liminf_{n\to\infty}\inf_{P\in\mathcal P}
     \Pr_{P^{(n)}}\{\theta(P)\in C_n\}\geq 1-\alpha,
 \qquad 0<\alpha<1,
\]
together with an informative diameter or expected-width requirement.
Coverage alone permits the vacuous whole parameter space. Testing questions
similarly impose both size control and power against a declared separated
alternative class.

\textbf{Binary causal baseline.}
When an entry explicitly invokes this baseline, one observes independent
copies of $O=(X,W,Y)$, with $W\in\{0,1\}$ and potential outcomes $Y(0),Y(1)$.
Assume consistency $Y=Y(W)$, no interference, conditional exchangeability
$(Y(0),Y(1))\perp W\mid X$, and overlap
$\epsilon\leq e(X)\leq1-\epsilon$ almost surely for a fixed
$0<\epsilon<1/2$, where
\[
 e(x)=\Pr(W=1\mid X=x),\qquad
 m_w(x)=\E[Y\mid W=w,X=x],\qquad
 \tau(x)=m_1(x)-m_0(x).
\]
These assumptions identify conditional mean effects almost everywhere
under the covariate law. Pointwise or uniform effects require appropriate
regularity and a specified support. Continuous treatment, longitudinal
data, adaptive experiments, and network interference require their own
assumptions in place of this baseline. Bounded outcomes, densities, and
smoothness are additional assumptions only where stated.

\textbf{Function classes.}
On $[0,1]^d$, $\mathcal H_d^s(L)$ is the isotropic Hölder ball of radius
$L>0$: put $k=\lceil s\rceil-1$ and $\gamma=s-k\in(0,1]$, bound derivatives
through order $k$, and require the order-$k$ derivatives to be
$\gamma$-Hölder with constant at most $L$. Thus integer orders use a
Lipschitz final derivative convention. The class
$\operatorname{BV}_d(L)$ consists of functions $f\in L^1((0,1)^d)$ whose
distributional derivative is a finite vector measure and for which
$\|f\|_{L^1}+|Df|((0,1)^d)\leq L$.
An $L^\infty$ bound or a particular pointwise version is imposed separately
when needed. Sparse and neural-network classes require a declared
dictionary or architecture and coefficient bounds.

An anisotropic H\"older class with declared block smoothness indices means
coordinatewise H\"older regularity in each block, uniformly in the other
blocks, with all corresponding derivative and H\"older bounds at most
the stated radius. Derivative orders and the integer-order convention in
each block are those just given. Mixed-derivative bounds are additional
restrictions only when explicitly stated.

\textbf{Decision and computational questions.}
Policy regret compares admissible feasible policies under the same law and
constraints. In computational entries, finite data and thresholds have
rational encodings; running time is measured in their total bit length.
Real-valued equilibrium search uses the explicitly declared algebraic or
FIXP encoding. An approximation ratio for maximization is written
$\operatorname{OPT}/\operatorname{ALG}$ and for minimization as
$\operatorname{ALG}/\operatorname{OPT}$, with zero optima and infeasible
instances handled locally. Historical approximation bounds are not
automatically current bounds.
% END ECONOMICS STATEMENT
\end{document}
